\section*{第 \ref{ch:g8:electric-current-effects} 章 --- 电流及其效应}

\begin{solution}{exo:g8:electric-current-effects:1}
热效应（烤面包机、电水壶、灯丝）；磁效应（
电磁铁、电动机、门铃）；化学效应（电镀、
电解水、给电池充电）。
\end{solution}

\begin{solution}{exo:g8:electric-current-effects:2}
热效应；磁效应；化学效应；热效应；磁效应。
\end{solution}

\begin{solution}{exo:g8:electric-current-effects:3}
电流能在导线周围产生磁性——这是电与磁之间的第一座
桥梁。电流方向反转，
小磁针的偏转方向也随之反转。
\end{solution}

\begin{solution}{exo:g8:electric-current-effects:4}
电磁铁可以随时断电、通电，还能随意改变磁极方向，
其磁性强弱也可以调节（增加匝数、增大
电流）。永久磁铁则以上这些都做不到。
\end{solution}

\begin{solution}{exo:g8:electric-current-effects:5}
若裸露的线圈匝直接相互接触，电流就会在绕组间短路，
而不再流经整个线圈——因此必须绝缘，迫使电流沿着
长长的磁化路径流动。另外，线圈本身对电流的阻碍很小：
若直接接在电池两端，几乎就是短路，线圈和电池都会发热——
这正是热效应的体现。
\end{solution}

\begin{solution}{exo:g8:electric-current-effects:6}
断开开关：电流消失，磁性消失，汽车被放下。
若是永久磁铁，磁性强到这种程度，将永远吸住汽车——
要松开它，就得动用机器把它硬拆下来。
\end{solution}

\begin{solution}{exo:g8:electric-current-effects:7}
放电时：储存的化学能转化为电能，驱动电流流动，
最终转化为光能、机械能和内能。充电时，
电能被迫反向流过电池，重建
化学能储备。
\end{solution}

\begin{solution}{exo:g8:electric-current-effects:8}
因为电流本身看不见，只能通过它的效应来察觉。
发热、磁针偏转、液体冒泡，正是电流仅有的标志——
而下一章的仪器会把磁针的偏转精确为可读的
数值。
\end{solution}

\begin{solution}{exo:g8:electric-current-effects:9}
按下按钮：电流接通，电磁铁吸住锤子，发出“叮”的一声——
但锤子摆动时会断开自身电源中的触点；磁性消失，弹簧把锤子弹回，
触点重新接通；电流恢复，再次吸住锤子——如此反复，
每秒可达多次。电磁铁无法持续吸住锤子，这是设计使然：吸合
总会切断它自己的电流来源——于是连续的颤响
\emph{即是}
铃声。
\end{solution}

\begin{solution}{exo:g8:electric-current-effects:10}
恒温器中脆弱的触点只承载线圈的小电流；线圈的磁性吸合
粗壮的触点，这些触点在另一个独立回路中承载加热器的大电流。
小电流控制大电流，两个电路彼此从不接触——
精巧的操作只需拨动杠杆，
而无须直接承受负载。
\end{solution}

\begin{solution}{exo:g8:electric-current-effects:11}
方法：将指南针平放在墙面上，远离铁质家具，
缓慢扫过墙面；磁针随可疑电器的开关而偏转之处，
即有通电导线经过。
偏转越明显，说明电流越近或越强。
需要注意：铁质水管、磁铁乃至地球本身，都会使磁针偏转，
所以只有偏转随开关的通断而同步变化，
才能证明存在电流。
\end{solution}

\begin{solution}{exo:g8:electric-current-effects:12}
电流使线圈变成磁铁；线圈与框架的异名磁极相互吸引，
同名磁极相互排斥，线圈便转向对齐的位置。
若不改变电流方向，线圈转到对齐位置便会停下——
转过半圈后抖动一下便归于静止，这只是门闭合器，而不是电动机。
在线圈恰好对齐时改变电流方向，线圈的两极随之对调：
原先的吸引变为排斥，对齐位置反而成了最不利的位置，
线圈必须不断追逐一个每转过半圈就躲开的目标——
这种永无止境的追逐，
也就是转动。
\end{solution}

\begin{solution}{pb:g8:electric-current-effects:1}
\textbf{1.} 三个演示实验\emph{串联}在同一个带开关的回路中：
同一股电流依次流经加热的导线、小磁针旁的导线和电解液——
一个开关，三种效应。
\textbf{2.} 串联规律：\emph{同一}电流无论流过粗导线还是细导线，大小都相同。
不同的只是路径——
细导线阻碍更大，因而发热更明显：
电流相同，受热各异。
\textbf{3.} 对调电池的正负极，\emph{并}把导线移到指南针的另一侧
（或将指南针调转）：两次反转相互抵消，
磁针的偏转方向仍指向东。
\textbf{4.} 水中两种成分以二比一的比例存在——
从一个电极得到的气体，
体积是另一个的两倍。
\textbf{5.} 去掉会自行断开的触点（即锤子摆动时断开的那个）：
电磁铁便会把毡制的锤子持续吸住，使其贴着铃而不发声——
参观者只能听到一声轻柔的“咚”，
此外再无声响。
\textbf{6.} 在废料箱旁，操作员断开线圈的开关：
磁性消失，铁块落下。铝块留在原处，因为
无论天然磁铁还是电流产生的磁铁，都只吸引铁这一类物质，
铝从来不受它们吸引。
\textbf{7.} 电流产生的磁铁，其磁极随电流方向而定：
改变其中一个，另一个也随之改变。
\textbf{8.} 电池靠\emph{消耗}自身储存的化学能来驱动电镀；
要重新储存化学能，需要从外部输入能量。
若电池能靠自己所做的功来补充自身能量，
便又回到了永动机的老问题——能量的收支
无法自己平衡自己。
\textbf{9.} 抓取：开关闭合，线圈磁化，夹住小玩具。
摆动：电流保持，抓力保持——磁效应正在起作用。
松开：开关在滑槽上方断开，磁性消失，
奖品落下。
\textbf{10.} 电磁铁依赖电流而存在：电池电压下降，电流减弱，
抓力随之减弱——可以随时通断，代价便是对电流的依赖。
（永久磁铁不会减弱，
也永远不会松开。）
\textbf{11.} 电流相同时，匝数越多，磁铁越强：
抓力越牢。另一种调节方式是电流本身的强弱——
即\emph{电流强度}，将在下一章中测量。
\textbf{12.} 例如：“热效应：保险丝，会率先忠实地熔断。
磁效应：抓娃娃机，听从指令抓取与释放。
化学效应：电镀槽，用电池的耐心
把勺子镀上银。”
\end{solution}
